\section{Struktur Logika Predikat: Term, Konstanta, dan Variabel}

Dalam tata bahasa, kalimat deklaratif sederhana umumnya memuat subjek dan predikat. Konsep serupa diadopsi dalam logika orde pertama. Komponen yang menunjuk objek yang dibicarakan disebut \logic{Term} \cite{rosen2019}.

\subsection{Definisi Term}

\begin{definisi}[Term]
Dalam matematika formal, \logic{Term} merujuk pada elemen atau entitas yang menjadi subjek bagi perekatan suatu pernyataan, sifat, atau relasi.
\end{definisi}

Term berperan sebagai \textbf{kata benda} (\textit{noun phrase}): ia bukan sifat atau relasi, dan tidak bernilai benar atau salah ($\text{True}/\text{False}$). Term menyatakan rujukan ke entitas.

Secara formal, bentuk term meliputi \textbf{tiga kategori}: konstanta, variabel, dan hasil aplikasi fungsi terhadap term.

\subsection{Konstanta (\textit{Constant})}

Konstanta adalah penamaan \textbf{spesifik} untuk suatu objek tertentu dalam semesta pembicaraan. Rujukannya tetap dan tidak bergantung pada konteks penafsiran bebas.

\begin{contoh}
\textbf{Contoh konstanta sebagai term:}
\begin{itemize}
    \item Nama entitas: ``Socrates'', ``Jakarta'' (lazim dinotasikan dengan huruf kecil, misalnya $s$, $k$).
    \item Bilangan matematika: $1$, $0$, $-45$, $3{,}14$, $\pi$.
    \item Objek konkret dalam narasi: ``buku ini'', ``laptop milik Budi''.
\end{itemize}
\end{contoh}

\subsection{Variabel (\textit{Variable})}

Variabel dipakai apabila subjek belum ditetapkan, sedang dicari, atau diwakili secara umum. Variabel dapat \textbf{disubstitusi} oleh konstanta dari semesta yang sesuai \cite{wiki_first_order_logic}.

\begin{contoh}
\textbf{Contoh variabel sebagai term:}
\begin{itemize}
    \item Simbol umum: $x$, $y$, $z$.
    \item Dalam kalimat \textit{``Ia adalah seorang presiden modern''}, kata ganti ``ia'' berperan seperti variabel $x$. Selama tidak diikat ke satu individu tertentu, kalimat tersebut belum menentukan proposisi tunggal tentang satu konstanta.
\end{itemize}
\end{contoh}

\subsection{Fungsi Pemetaan (\textit{Function})}

Fungsi (dalam bahasa logika) memetakan satu atau lebih term ke sebuah \textbf{term baru}. Nilai fungsi tetap merupakan term, bukan nilai kebenaran.

\begin{contoh}
\textbf{Contoh fungsi sebagai term:}
\begin{itemize}
    \item Ekspresi aritmetika: $f(x) \equiv x^2 + 1$. Untuk $x=3$, diperoleh term bilangan $f(3)=10$.
    \item Fungsi deskriptif: $ayah(x)$ bermakna ``ayah kandung dari $x$''. Untuk konstanta tertentu, $ayah(c)$ merujuk pada individu tertentu (konstanta) dalam semesta manusia.
    \item Fungsi terstruktur: $dekan(F)$ bermakna ``individu yang menjabat dekan pada fakultas $F$''. Dalam suatu semesta pembicaraan dan periode waktu tertentu, ekspresi ini dapat disamakan dengan satu konstanta individu (misalnya $d$) tanpa menamai orang secara spesifik di teks.
\end{itemize}
\end{contoh}
